\documentclass[a4paper,12pt]{article}
\usepackage[T1]{fontenc}
\usepackage{lmodern}
\usepackage[margin=25mm]{geometry}
\usepackage{parskip}
\usepackage{fancyhdr}
\setlength{\parindent}{0pt}
\setlength{\parskip}{6pt}
\linespread{1.06}
\pagestyle{fancy}
\fancyhf{}
\renewcommand{\headrulewidth}{0pt}
\fancyfoot[C]{\small\thepage}
\newcommand{\reflectionheading}[1]{\par\addvspace{9pt}{\large\bfseries #1}\par\nopagebreak}
\begin{document}

\begin{center}
{\LARGE\bfseries Individual Self-Reflection Report}\\[5pt]
{\small II2202 -- Research Methodology and Scientific Writing}
\end{center}

Lorenzo Deflorian\\
Group number: 4\\
Project: Transparent Multi-Hop Routing for Raft Under Partial Network Partitions

\reflectionheading{What I did in the project}
In our Research Plan, I was responsible for leading the Raft, client-workload, and measurement components, together with the analysis of progress, recovery, and elections for RQ1. My work followed this division, although some tasks overlapped as we combined the different parts.

I worked on the experimental setup, checked the collected results, and helped prepare the figures and report. I also contributed to the AWS setup used to run the experiments and to revising the report after the opposition feedback.

Some practical work received little attention in the final report because we needed to focus on the research questions and results. This included debugging the setup, checking logs, automating runs, and making sure the collected data was complete and usable.

\reflectionheading{Teamwork}
Overall, I think we worked well together. We had separate areas of responsibility, which helped us make progress, but we also discussed the setup and results together. Our tasks were different, but I felt that we contributed roughly equally in terms of time and effort.

Combining the components required coordination and some joint debugging. There were no major teamwork issues that I think the examiner needs to be aware of.

\reflectionheading{What I enjoyed}
I most enjoyed turning a question about distributed systems into an experiment that we could actually run. It was interesting to see how failures in individual network links affected Raft and whether forwarding could help.

I am most proud of getting the complete experiment working and collecting results from 900 trials. Seeing the difference between direct communication and forwarding made the project feel concrete.

\newpage
\reflectionheading{What I learned}
I developed a better understanding of Raft, especially how leader elections and communication with a majority affect progress. I also learned more about designing repeatable experiments, collecting measurements, and interpreting results carefully.

One important lesson was that a failed client request does not necessarily mean that the consensus protocol has stopped working. We needed to distinguish problems with client access from problems with the cluster itself.

The project also helped me practise scientific writing: connecting results to research questions, explaining limitations, and keeping claims within what the evidence supports.

\reflectionheading{What I found difficult}
The hardest part was making sure that our measurements represented what we intended to study. Recovery times could include setup work, and client placement could affect whether requests succeeded.

Inspecting logs and checking the experiment sequence helped us understand these effects. If I repeated the project, I would spend more time validating the measurements during the pilot stage. This would make the later analysis easier and reduce uncertainty when interpreting the results.

\end{document}
